\begin{abstract}
Small language models deployed on a device could keep improving by learning from their own work, but only if something decides which experiences are worth learning from. We study this in a ``creature'': a frozen Qwen3 model (0.6B or 1.7B) with a small LoRA adapter that works by day and, each night, trains on the episodes judged good, replaying recent lessons. Over lives of up to 60 days we compare judges modelled on published approaches: the model's own judgment (as in self-rewarding models), its reflection on the user's reaction, the user's stated verdict read without error, a separate frozen ``feedback sense'' model that reads the reaction, and an oracle that knows the true outcome (as in training on verified successes).

On WebShop, learning from its own judgment decayed in every life, while the oracle improved and then held steady for 60 days. A first feedback sense (v2) let a few polite give-ups through as successes, and 3 of 5 lives collapsed after learning from them on the first night. A second (v3), built for discernment, which first asks whether the user gave a clear verdict at all, cut the reactions it misread as approval from 5.0\% to 1.2\%, and 4 of 5 lives reached near-oracle scores within 14 days. They then declined slowly, although their lessons stayed clean: about 2 a day against the oracle's 20, with only 2-10\% failures. Across 97 lives, the share of failed attempts among a life's lessons predicted where it ended (Spearman $\rho$ = -0.81): a handful of correct lessons was enough to improve, and a handful of wrong ones enough to collapse.

On $\tau^2$-bench, where a simulated customer cannot see whether its request was actually carried out, 80\% of the lessons sense v3 accepted were failed tasks, and all three 1.7B sense v3 lives collapsed within 30 days. Even oracle lessons did not raise overall scores: they came mostly from the easiest domain, which improved while the other learned domain and a never-practised one declined.

Learning from users works only when the judge keeps failed attempts out of the lessons, which needs users who can see the outcome; and even correct lessons must be balanced across skills to avoid forgetting.
\end{abstract}
